\chapter{Baseline Investigation \& The SOTA Repair Framework}
\label{ch:sota-repair}

A comparison against a published state-of-the-art model is only as meaningful as the
baseline it runs. This chapter documents what happened when Brain-TokenGT --- a MICCAI 2023
spatiotemporal graph transformer for longitudinal connectome embedding --- was taken from
its open-source release and run on the DELCODE cohort. The short version is that the
baseline required three successive interventions before it produced a number at all, and
that the number it eventually produced is \emph{not reproducible at a fixed random seed}.
That second finding is the more important one, and it generalises well beyond this
particular model.

%==============================================================================
\section{Brain-TokenGT as a Published Benchmark}
\label{sec:btgt-what}

Brain-TokenGT was selected as the comparison baseline because it is the closest published
match to this thesis's problem statement: it consumes a \emph{sequence} of functional
connectomes rather than a snapshot, it was evaluated on dementia-conversion prediction in
MCI among its three tasks, and a reference implementation is publicly available.

Architecturally it comprises two modules. Graph Invariant and Variant Embedding (GIVE)
produces node and spatio-temporal edge embeddings, which are tokenised. The Brain Informed
Graph Transformer Readout (BIGTR) then augments those tokens with trainable type
identifiers and non-trainable node identifiers and passes them to a standard transformer
encoder. The temporal mechanism inside GIVE derives from EvolveGCN-H, whose graph
convolutional recurrent unit (GRCU) evolves the graph convolution's \emph{weight matrix}
through a GRU rather than propagating a node hidden state.

The tokenisation is where the capacity goes. With $K = 8$ nodes per ROI clique, one visit
yields roughly 460 tokens; across a three-visit sequence the transformer attends over
approximately $1{,}381$ spatio-temporal tokens per subject, through $603{,}849$ trainable
parameters. For context, the cross-validation folds used here contain roughly 76 training
subjects each.

%==============================================================================
\section{Uncovering Baseline Failure Modes}
\label{sec:btgt-failures}

\subsection{The as-released recurrent core is inert}

In the released configuration, the flag \texttt{train\_give} is set to \texttt{false}. The
consequence is not a reduced learning rate on the recurrent component --- it is that the
GRCU weights are frozen at their random Gaussian initialisation and never registered with
the optimiser at all. The model has a temporal module in the architectural sense and no
learned trajectory in any operational sense.

The behaviour on DELCODE is the signature of a model that never learned a decision
boundary: test AUC $0.5779$, sensitivity $1.000$, specificity $0.000$ --- every test
subject predicted as a converter --- with cross-validation AUC $0.5538 \pm 0.1258$.

\risk{Reporting a win over the as-released configuration would be a null comparison.}{%
A model whose temporal core was never trained is not a baseline; it is an untrained
control. Any claim of the form ``our model beats Brain-TokenGT'' resting on the
\texttt{train\_give=false} number is indefensible, and this thesis does not make it. The
as-released number is reported only to document why repair was necessary.}

\subsection{Enabling the recurrent core diverges}

Setting \texttt{train\_give=true} and allowing end-to-end gradient updates to the recurrent
weights produces numerical collapse rather than learning. Gradient norms in the GRCU
parameters explode, non-finite scores reach \texttt{TopK.forward}, and training crashes
within approximately ten epochs. Two consecutive attempts failed outright with a non-zero
exit code.

This is itself a finding worth stating plainly: standard EvolveGCN-H is intrinsically
unstable when trained end-to-end on clinical cohorts of $N < 100$.

%==============================================================================
\section{The Stabilisation Protocol --- and What It Does Not Achieve}
\label{sec:btgt-repair}

The intervention that produced completing runs was a parameter-group decoupling: the
newly-unfrozen recurrent parameters were placed in their own optimiser group with a damped
learning rate (\texttt{give\_lr\_scale} $= 0.1$) and weight decay
(\texttt{give\_weight\_decay} $= 10^{-3}$), leaving the remaining parameters untouched.
Under this configuration the model completes training reliably --- $5/5$ runs in the
initial batch, and $19/19$ across the full sweep reported below.

It is tempting, and it was initially tempting in this project, to call that arm
``stabilised'' and move on. That label is wrong, and the distinction it obscures is the
subject of the rest of this chapter.

\begin{quote}
  \textbf{``Stabilised'' as established here means \emph{crashes less often}. It does not
  mean \emph{reproducible}.}
\end{quote}

%==============================================================================
\section{The Reproducibility Audit}
\label{sec:btgt-reproducibility}

\subsection{Design}

The audit asks a question that the brain-graph literature does not appear to ask anywhere
(Section~\ref{sec:gap-determinism}): \emph{if the identical configuration is run twice at
the identical seed, does it produce the identical result?} Answering it requires repeat
runs at a fixed seed, not a multi-seed sweep. Nineteen runs of the stabilised configuration
were collected --- seven at seed 42, four each at seeds 43, 44, and 45 --- all completing
successfully, all sharing the same code state.

\subsection{Result}

\begin{table}[htbp]
  \centering
  \small
  \begin{tabular}{@{}rrrrl@{}}
    \toprule
    \textbf{Seed} & \textbf{$n$ runs} & \textbf{Mean test AUC} & \textbf{Within-seed SD} & \textbf{Range} \\
    \midrule
    42 & 7 & 0.5779 & 0.1167 & 0.3571 -- 0.7078 \\
    43 & 4 & 0.6802 & 0.1024 & 0.5779 -- 0.8182 \\
    44 & 4 & 0.5227 & 0.0849 & 0.4156 -- 0.6234 \\
    45 & 4 & 0.6477 & 0.1002 & 0.5195 -- 0.7338 \\
    \midrule
    \multicolumn{2}{@{}l}{\textbf{All 19 runs}} & \textbf{0.6025} & 0.1123 (pooled) & 0.3571 -- 0.8182 \\
    \bottomrule
  \end{tabular}
  \caption{Brain-TokenGT (stabilised) held-out test AUC across 19 completed runs of the
    identical configuration. Seed 42 alone spans an AUC range of $0.35$ --- from well below
    chance to the best result obtained at any seed.}
  \label{tab:btgt-19runs}
\end{table}

The decomposition is the headline:

\begin{align*}
  \text{mean within-seed SD} &= 0.1011 \\
  \text{between-seed-mean SD} &= 0.0706
\end{align*}

\textbf{Run-to-run variance exceeds seed-to-seed variance, at every seed, not merely at
seed 42.} A ``four-seed'' summary of this model is therefore four draws from run-to-run
noise wearing the label of a seed effect. The four seeds are not four measurements of four
conditions; they are four samples from one distribution.

\subsection{Mechanism}

The cause is not mysterious and is documented in the experiment registry itself: only
\texttt{cudnn.deterministic} is set by the shared seeding helper. The codebase never calls
\texttt{torch.use\_deterministic\_algorithms(True)} and never sets
\texttt{CUBLAS\_WORKSPACE\_CONFIG}, so operations whose backward pass accumulates through
non-deterministic GPU atomics remain non-deterministic regardless of seeding. Brain-TokenGT
reaches two such operations that the recurrent trajectory models of Chapter~\ref{ch:methodology}
do not: \texttt{TopKPooling} in the transformer module, and the index-gather over
\texttt{node\_embs} in the GRCU. Floating-point reduction order varies between runs, and on
a model already sitting near an optimisation cliff, that is sufficient to move the outcome
by a third of an AUC point.

\subsection{The contrast: this thesis's own pipeline is bit-reproducible}

The same experiment harness, the same seeding helper, the same GPU. The \texttt{none}-arm
GELSTM configuration was run twice at seed 42 and produced \textbf{byte-identical} results
(test AUC $0.7607$, CV AUC $0.9010$ both times). The non-determinism is therefore a
property of the audited baseline's operations, not of the shared infrastructure.

\edge{Reproducibility itself is a reportable result of this comparison.}{%
``Our pipeline re-runs byte-identically at a fixed seed; the published baseline spans
$0.35$ AUC at a fixed seed'' is a stronger, more defensible, and more useful statement than
any accuracy margin between the two --- and unlike an accuracy margin on $N=25$ test
subjects, it does not depend on cohort size.}

\risk{The determinism of this thesis's \emph{frozen-encoder} arm is asserted from the
\texttt{none} arm, not measured directly.}{%
Bit-reproducibility was verified for the \texttt{none} arm, which builds no encoder and so
executes no scatter backward. The \texttt{pretrained\_frozen} arm used in the matched-cohort
comparison has exactly one run per seed, so its own within-seed variance is unmeasured. The
frozen encoder is not trained, so no scatter \emph{backward} runs through it --- but this is
an argument, not a measurement, and it should be replaced by a repeat run before the claim
is generalised to the whole pipeline.}

%==============================================================================
\section{Consequences for the Statistical Comparison}
\label{sec:btgt-stats}

Two paired significance tests were computed in an earlier iteration of the matched-cohort
comparison. Both are invalid, for distinct reasons, and both are withdrawn.

\begin{enumerate}
  \item \textbf{Held-out test AUC, paired $t$-test, $n = 4$ seeds ($p = 0.0299$).} This
    treats seed as the sole source of variance. Section~\ref{sec:btgt-reproducibility}
    shows the within-seed term is the \emph{larger} one, so the quantity the test assumes
    it is measuring is not the quantity that varies. The $p$-value is uninterpretable.
  \item \textbf{Cross-validation AUC, paired $t$-test, $n = 20$ folds
    ($p = 1.59 \times 10^{-6}$).} This treats the five folds within a seed as five
    independent paired units. Folds within a seed share training data by construction and
    are strongly correlated; the effective sample size is nearer 4 than 20. No correction
    rescues a test whose independence assumption fails this directly, and the reported
    $p$-value is inflated by roughly the fold-correlation factor.
\end{enumerate}

\risk{An earlier version of this comparison quoted one run per seed out of seven available
on disk.}{%
The comparison notebook hard-coded a single run path per seed, so seed 42 was represented
by a run scoring $0.6104$ while six other completed seed-42 runs --- spanning $0.3571$ to
$0.7078$ --- sat unexamined in the same output tree. This was selection by accident rather
than by intent, but the effect is identical: a reported mean of $0.6672$ against a true
grand mean of $0.6025$. The reporting path now globs every run directory per experiment
identifier so that a stale or unrepresentative path cannot silently become ``the'' result.}

\paragraph{What survives.} The \emph{direction} of the comparison is robust and does not
depend on any of the withdrawn statistics. The GELSTM matched-cohort arm attains
$0.8782 \pm 0.0256$ across four seeds; the Brain-TokenGT distribution over 19 runs has a
grand mean of $0.6025$ and a maximum of $0.8182$. The best Brain-TokenGT run ever observed
is below the GELSTM mean. That comparison is properly reported as two distributions and an
effect size, never as a $p$-value.

%==============================================================================
\section{Methodological Fairness \& Discrepancy Alignment}
\label{sec:btgt-fairness}

Three asymmetries remain between the two arms and are stated rather than minimised.

\subsection{Cohort windowing --- resolved}

Brain-TokenGT requires $2 \le T \le 3$ visits, whereas the full DELCODE longitudinal cohort
spans $1 \le T \le 6$. Comparing across that difference would confound architecture with
sequence length and with a shifted task difficulty: single-visit subjects test static
separability, not trajectory modelling. The comparison reported here therefore restricts
\emph{both} models to the identical $2 \le T \le 3$ window ($N_{\text{CV}} = 95$: 37
converters, 58 stable MCI; $N_{\text{test}} = 25$: 11 converters, 14 stable MCI), with
subject-identifier and fold-membership parity verified programmatically rather than assumed
from a shared seed.

\subsection{Capacity and tokenisation --- not resolved, and not resolvable}

\begin{table}[htbp]
  \centering
  \small
  \begin{tabular}{@{}lrrr@{}}
    \toprule
    \textbf{Architecture} & \textbf{Trainable params} & \textbf{Temporal-core input} & \textbf{CV $\rightarrow$ test gap} \\
    \midrule
    GELSTM (\texttt{none})   & $31{,}169$  & 1 vector / visit       & $+0.060$ \\
    GELSTM (frozen GAAE)     & $520{,}905$ & 1 vector / visit       & $+0.140$ \\
    Brain-TokenGT (stabilised) & $603{,}849$ & ${\sim}1{,}381$ tokens & $-0.198$ \\
    \bottomrule
  \end{tabular}
  \caption{Parameter and representation asymmetry. The sign of the CV-to-test gap is the
    interpretable quantity: Brain-TokenGT's high validation performance does not transfer,
    which is the classic small-sample overfitting signature.}
  \label{tab:capacity}
\end{table}

This asymmetry cannot be engineered away, and arguably should not be: it \emph{is} the
finding. A transformer attending over ${\sim}1{,}381$ tokens per subject, trained on ${\sim}76$
subjects per fold, is being asked to do something the sample size does not support. Its
validation AUC of $0.8142 \pm 0.0838$ collapsing to a test AUC of $0.6156 \pm 0.0705$ is
what that looks like from the outside.

\subsection{Tuning budget}

The recurrent trajectory models were developed iteratively across many experiments, while
Brain-TokenGT was run at its author-recommended defaults with only the minimal
learning-rate and weight-decay adjustment required to prevent divergence. Under the
fairness criterion adopted for this project --- that all models be \emph{decently}
optimised, meaning no collapsed learning rate and no training continued past obvious
overfitting, without requiring exhaustive search --- this is acceptable. It is not
equivalent tuning, and no claim in this thesis depends on it being equivalent.

\risk{Brain-TokenGT received no search over its own structural hyperparameters.}{%
Token count $K$ and transformer depth $L$ were left at published defaults. Given that the
capacity/sample-size mismatch of Table~\ref{tab:capacity} is central to the interpretation,
a reviewer may reasonably ask whether a smaller $K$ would have closed the gap. The honest
answer is that it was not tested, and that a fair version of that test would require an
equivalent structural search on the recurrent side.}

\subsection{Temporal representation}

The recurrent models condition their gate updates on the actual inter-visit interval in
months, distinguishing a six-month follow-up from a twenty-four-month one. Brain-TokenGT
treats visits as uniform discrete steps. This is a genuine architectural difference in
favour of the models developed here rather than an unfairness in the setup --- but it does
mean the comparison measures ``$\Delta t$-aware recurrence versus $\Delta t$-agnostic
attention,'' not ``recurrence versus attention'' in isolation.
